% Cote 119: offprints, reports, courses, CV and letters (typed and printed part): archive cote 119, site page 193 (PDF page 194).
% typescript. Draft compared: none.
% AI-transcribed by Claude Opus 5.5 (claude-opus-5-5, adaptive thinking, effort high), Message Batches API, batch r119_full (msgbatch_01DwZNm5Y7YsKBWyAN4sFs8w), 2026-11-13; not yet checked by eye.
\scanpage{119/193}
UNIVERSITE DES SCIENCES ET \hfill Année Universitaire 1978-79

TECHNIQUES DU LANGUEDOC

MATHEMATIQUES

EN GUISE DE PROGRAMME

pour le cours de C4 de A. Grothendieck ``Introduction à la recherche'' 1978/79

Quand une curiosité intense anime une recherche, nous avançons comme
portés par des ailes impatientes. Ne sommes-nous alors téméraire esquif aux
voiles tendues qui avidement laboure l'inépuisable océan ? Oui, de toutes
parts sommes-nous entourés de brumes mouvantes, sans cesse prenant corps et
s'éclairant sous les yeux qui les fouillent, sans cesse se dérobant pour mieux
nous provoquer à les pénétrer ! Et nous exultons en le mystère de toute nou-%
velle énigme entrevue, par notre regard pressant dépouillée de ses voiles de
brume, pour être connue féconde en mystères nouveaux...

L'ardente curiosité seule est créatrice, elle nous porte droit au
coeur même de l'Inconnu. N'est-ce pas Elle notre seul et véritable héritage,
déposé en chacun de nous dès avant que nous fussions enfantés ? Graine imper-%
ceptible, dont pourtant nait la Fleur aux mille pétales comme l'Arbre à la
ramure innombrable... Il n'y a rien qui ne naisse d'Elle. Et pour peu que nous
la laissions s'épanouir en nous, il n'y a rien que ne puisse enfanter notre
seule Soif de connaître. Elle seule nous donne des ailes, Elle seule anime
l'élan qui nous porte au coeur des choses. Où elle n'est, il n'y a Création,
ni Amour.
Quand cette soif est absente, quel sens reste-t-il à notre vie ?
Quel sens a un travail où il n'y a création, ni amour ? Que reste-t-il donc,
quand il ne semble plus y avoir trace de l'enfant en nous qui joue et qui
interroge ? Quel est l'avenir d'un monde qui laisse périr son unique héritage ?

Les trois dernières années, j'ai enseigné comme un aveugle peindrait.
Je parlais de choses que je découvrais à mesure, à des personnes venues m'écou-%
ter par quelques étrange obligation. Certes, les choses vues et dites étaient
si tangibles et si simples qu'un enfant curieux les pouvait découvrir avec moi
comme compagnon de jeu - et je parlais comme j'aurais parlé à cet enfant, ou
à moi-même. Et porté par ce dialogue imaginaire, je restais aveugle au fait que
je monologuais, devant des écoliers appliqués à prendre les notes d'un cours
qui ne les concernait pas. Seul l'examen les concernait. Les choses dites
avaient beau être enfantines et vivantes - c'est comme autant d'objets hété-%
roclites et morts que pêle-mêle elles s'entassaient pesamment en des esprits
inertes - frappés de paralysie.

L'indifférence est à jamais impuissante à étreindre, fût-ce l'évi-%
dence que l'enfant reconnaît en jouant. Oui, par quelque effort qu'elle s'éver-%
tue pour parvenir à ses fins, l'indifférence reste impuissante. Quand il n'est
mû par la joie, l'effort souvent certes débouche sur l'angoisse, jamais sur
une compréhension. Où il n'y a compréhension, peut-il y avoir maîtrise ?

A mon insu, prisonnier inconscient des charmes d'un solitaire
voyage de découvertes, je n'ai fait que perpétuer en des écoliers sans voix
les vieilles angoisses, les vieilles impuissances. Quelques notes de fin
d'année, griffonnées d'une main lasse sur des copies d'examen écrites sans

\hfill .../...

\mspage{\hand{193}}
